\documentclass[10.5pt,a4paper]{article}

\usepackage[T1]{fontenc}
\usepackage[utf8]{inputenc}
\IfFileExists{lmodern.sty}{\usepackage{lmodern}\usepackage{microtype}}{}
\usepackage[english]{babel}
\usepackage[top=1.3cm,bottom=1.2cm,left=2.0cm,right=2.0cm]{geometry}

\setlength{\parskip}{0.35em}
\setlength{\parindent}{0pt}

\begin{document}
\pagestyle{empty}

%% ============================================================
%% Cover letter [First] [Last], Bain & Company Zurich
%% Associate Consultant. English, one page.
%% No em-dashes; only normal hyphens in compounds.
%% ============================================================

%% Sender
\begin{flushleft}
[First] [Last]\\
Zürich, Switzerland\\
{[Phone]}\\
{[your.email@example.com]}
\end{flushleft}

%% Recipient
\begin{flushright}
Bain \& Company Switzerland\\
Recruiting Team\\
Sihlporte 3\\
8001 Zürich
\end{flushright}

Zürich, 28 August 2026

\vspace{0.15em}

\textbf{Application: Associate Consultant}

Dear Recruiting Team,

Two weeks of fieldwork at Zurich Airport produced roughly 100,000 short passenger interviews, and when I began the analysis one transport mode had apparently doubled since the previous year. The easy call was to log it as a data error and move on. Instead I wrote down every explanation I could think of, ranked them by likelihood, and worked through them in order: interviewer effects, coding and digitisation errors, sample composition, then signage, construction, passenger routing. None of them held. So I went looking for independent evidence and found that the airport's escalators had automated passenger counters, a dataset nobody on the team knew existed. They confirmed the volume was real. Rolling stock and timetable data from SBB explained why: not more trains, but bigger ones. The finding fed into Flughafen Zürich's multimodal planning and infrastructure prioritisation.

I lead with that because it is the clearest answer I have to why I want to be an Associate Consultant at Bain. Hypothesis-led thinking is not a method I would be learning here. It is how I was trained to work, and how I have worked since.

The training came from ETH Zürich, where I took a Master's in Comparative and International Studies focused on causal inference and experimental design. The discipline was to start from a claim rather than from the data: write down what you believe is happening, what the world should look like if you are right, and what evidence would prove you wrong. My thesis applied that to short-term rental regulation, pairing qualitative coding of 146 regulatory documents from 133 US cities with a difference-in-differences estimate of the effect on housing costs. A separate ETH project, survey experiments across eight European cities in eight languages, became a co-authored article in \textit{Nature Sustainability}.

Since then I have worked both sides of a line most analysts only see from one side. At Sotomo I owned analytical workstreams end to end for corporate and public clients: a national health-behaviour study for CSS, and a study of why more than 10,000 Zurich residents did or did not use the city's digital services, each running from the first client conversation through to the recommendation. At Interface I did the same for policy, evaluating programmes for federal, cantonal, and municipal clients, including one for the Swiss Parliamentary Control of the Administration. What that range taught me is that the analysis is usually the easy half. The hard half is knowing the difference between something a commercial client can act on next quarter and something a public body has to be able to defend in public.

I also build the tools rather than only using them. At Sotomo I developed the internal R packages the team adopted as its standard analysis infrastructure, and I integrated generative AI into my R workflows to summarise qualitative input, generate first-pass hypothesis sets, and flag anomalies. That shortened timelines without moving judgement off the human. Your posting treats generative AI as part of how the work gets done rather than as a novelty, which is already how I work.

What I want from the Associate Consultant role is the half I have only seen from the outside: carrying a recommendation through to a decision that actually gets made, on a client's timeline, with the people who have to live with it. Bain does that for a corporate client and a public institution out of the same office, and that range is exactly what I have been building toward.

I work in English and German at C2 and I am based in Zürich. I would welcome the chance to talk it through.

\vspace{0.2em}

Kind regards

\vspace{0.7em}

[First] [Last]

\end{document}
